\documentclass[12pt]{article}

% Ustawienia strony
\usepackage{lipsum}
\usepackage[margin=1.5cm]{geometry}
\usepackage{graphicx}
\usepackage{enumitem}
\usepackage{siunitx}
\usepackage{amsmath}
\usepackage{cancel}
\usepackage{makecell}
\usepackage{mathtools}
\usepackage{floatrow}
\usepackage{caption}
\usepackage[T1]{fontenc}
\usepackage{fancyvrb}
\usepackage[ngerman]{babel}
\usepackage{subcaption}
\usepackage{graphicx}
\usepackage{blindtext}
\usepackage{hyperref}
\usepackage[dvipsnames]{xcolor}

\hypersetup{
	colorlinks=true,
	linkcolor=blue,
	filecolor=magenta,      
	urlcolor=cyan,
	pdfpagemode=FullScreen,
}

\urlstyle{same}

\fvset{tabsize=4}
\setlist{leftmargin=15pt,labelindent=15pt}
\setlist[enumerate]{noitemsep,wide=0pt, leftmargin=*}
{\setlength{\parindent}{0pt}
	
	% Ustawienia stopki
	\usepackage{lastpage}
	\usepackage{fancyhdr}
	\pagestyle{fancy}
	\fancyhf{} % sets both header and footer to nothing
	\renewcommand{\headrulewidth}{0pt}
	\cfoot{\thepage\ z \pageref*{LastPage}}
	
\begin{document}

\section*{Podstawowe dane}
\begin{figure}[htp]
    \begin{minipage}[t]{0.3\textwidth}
        \textbf{Imie i nazwisko} \\
        \\
        \textbf{Data} \\
        \textbf{Tytuł} \\
        \textbf{Numer projektu}
    \end{minipage}%
    \begin{minipage}[t]{0.7\textwidth}
        Wojciech Kaźmierczak \\
        Robert Mesek \\
        01.01.2024 \\
        Projekt z Algorytmów Geometrycznych \\
        7
    \end{minipage}%
\end{figure}

\section*{Dane techniczne}
\begin{figure}[htp]
    \begin{minipage}[t]{0.3\textwidth}
        \textbf{Procesor} \\
        \textbf{Język} \\
        \textbf{Biblioteki} \\
        \\
        \textbf{System operacyjny} \\
        \textbf{Środowisko}

    \end{minipage}%
    \begin{minipage}[t]{0.7\textwidth}
        Intel i7-1165G7 x86\_64 \\
        Python 3.11.5 \\
        numpy 1.25.2, pandas 2.0.3, matplotlib 3.7.2, \\
        notebook 6.5.4, sortedcontainers 2.4.0 \\
        WSL Ubuntu 22.04 (MS Windows 10 22H2) \\
        conda 23.11.0
    \end{minipage}%
\end{figure}

\section*{Opis projektu}
Należy zaimplementować algorytmy konstrukcji triangulacji dla dowolnego wielokąta prostego dwoma metodami.
W obu przypadkach triangulacja ma się opierać jedynie na wierzchołkach wielokąta (bez dodawania wierzchołków wewnątrz wielokąta).

Pierwsza metoda to podział wielokąta na wielokąty monotoniczne (zgodnie z algorytmem podanym na wykładzie), a następnie triangulacja wynikłych z podziału wielokątów monotonicznych także metodą podaną na wykładzie (i realizowaną na ćwiczeniach).

Drugi algorytm to triangulacja Delaunay’a z ograniczeniami (z procedurami odzyskania krawędzi i usunięcia trójkątów zewnętrznych). W tej metodzie wybrać odpowiedni sposób poszukiwania trójkąta w istniejącej triangulacji (i go omówić).

\newpage

\section*{Triangulacja polegająca na podziale wielokątów}

\newpage

\section*{Triangulacja Delaunay'a z ograniczeniami}

\end{document}
